%!TEX root = TheoreticalPhysicsC.tex

\section{Lecture: Hilbert Space, Bases, Operators}

Let $V$ be a complex vector space. In quantum mechanics the following notation and conventions are used universally. From now on we shall write every vector $v\in V$ in the so-called ``ket'' notation:
\begin{equation}
	v \mapsto |v\rangle.  
\end{equation} 
This notation means nothing more than: ``$|v\rangle$ stands for a vector in a vector space $V$''. Why we have introduced this unusual notation will become clear shortly.

The effort pays off: the states of a quantum mechanical system are modelled by vectors in a Hilbert space. The real hurdle at the beginning is not the physics, but the notation.

We shall carry out a few calculations to make ourselves familiar with this notation.
\begin{example}
	Suppose $V\equiv \mathbb{C}^2$. (In QM we shall later interpret this Hilbert space as the \emph{state space} of a qubit.) The dimension of $V$ is $2$. Two example vectors are
	\begin{equation}
		|e_0\rangle \equiv \begin{pmatrix}
			1 \\ 0
		\end{pmatrix} \quad\text{and}\quad |e_1\rangle \equiv \begin{pmatrix}
			0 \\ 1
		\end{pmatrix}.
	\end{equation}
	These two vectors can be added, subtracted, and so on. For example, we can form the following vectors
	\begin{equation}
		|e_+\rangle \equiv \frac{1}{\sqrt{2}}|e_0\rangle  + \frac{1}{\sqrt{2}}|e_1\rangle \equiv \frac{1}{\sqrt{2}}\begin{pmatrix}
			1 \\ 0
		\end{pmatrix} + \frac{1}{\sqrt{2}}\begin{pmatrix}
			0 \\ 1
		\end{pmatrix} = \frac{1}{\sqrt{2}}\begin{pmatrix}
			1 \\ 1
		\end{pmatrix}
	\end{equation}
	and
	\begin{equation}
		|e_-\rangle \equiv \frac{1}{\sqrt{2}}|e_0\rangle  - \frac{1}{\sqrt{2}}|e_1\rangle \equiv \frac{1}{\sqrt{2}}\begin{pmatrix}
			1 \\ 0
		\end{pmatrix} + \frac{1}{\sqrt{2}}\begin{pmatrix}
			0 \\ -1
		\end{pmatrix} = \frac{1}{\sqrt{2}}\begin{pmatrix}
			1 \\ -1
		\end{pmatrix}.
	\end{equation}
\end{example}

Note that it makes no difference whether we write $|v+w\rangle$ or $|v\rangle+|w\rangle$, that is to say 
\begin{equation}
	|v+w\rangle \equiv |v\rangle + |w\rangle.
\end{equation}


The ket notation is generalised by introducing so-called ``bra'' vectors, which in the finite-dimensional case stand for the row vectors that result from the dagger operation: suppose 
\begin{equation*}
|v\rangle \equiv \begin{pmatrix}
	v_1 \\ v_2 \\ \vdots\\ v_n
\end{pmatrix}
\end{equation*}
were written in ket notation as $|v\rangle$, then $\langle v|$ would stand for the row vector
\begin{equation}
	\begin{pmatrix}
		v_1^* & v_2^* & \cdots & v_n^* 
	\end{pmatrix}
\end{equation}
instead. We define the result of the dagger operation applied to a ket vector as follows:
\begin{equation}
	(|v\rangle)^\dagger \equiv \langle v|.
\end{equation}

With this notation the inner product of two vectors $|v\rangle$ and $|w\rangle$ can be expressed very efficiently as follows:
\begin{equation}
	\langle v|w\rangle \equiv (|v\rangle)^\dag |w\rangle.
\end{equation}
(The name ``bra-ket'' goes back to this definition.)

\begin{historical}
	The notation is due to Paul Dirac. The word \emph{bracket} is split in two: on the left stands the ``bra'', on the right the ``ket'', and the ``c'' has gone missing.
\end{historical}

\begin{example}
	In the case $V\equiv \mathbb{C}^2$ the inner product of $|e_0\rangle$ and $|e_+\rangle$ is calculated as follows:
	\begin{equation}
		\langle e_0|e_+\rangle \equiv (|e_0\rangle)^\dag|e_+\rangle \equiv \begin{pmatrix}
		1 & 0
		\end{pmatrix} \begin{pmatrix}
			\frac{1}{\sqrt{2}} \\ \frac{1}{\sqrt{2}}
		\end{pmatrix} = \frac{1}{\sqrt{2}}.
	\end{equation}
\end{example}

It is important to note that a linear combination of two ket vectors behaves as follows under the dagger operation:
\begin{equation}
	(\alpha |v\rangle +\beta |w\rangle)^\dag = \alpha^* \langle v| + \beta^*\langle w|.
\end{equation}
It is exactly this rule that makes the notation useful: we work out inner products without ever writing the vectors down as columns -- in infinite-dimensional spaces the only way to calculate at all.

\begin{example}
	The inner product $\langle v|w\rangle$ of
	\begin{equation}
		|v\rangle \equiv \begin{pmatrix}
			1 \\ i
		\end{pmatrix} \quad \text{and} \quad |w\rangle \equiv \begin{pmatrix}
			1 \\ 1
		\end{pmatrix}
	\end{equation}
	can be calculated as follows. First, note that 
	\begin{equation}
		|v\rangle = |e_0\rangle + i|e_1\rangle \quad \text{and} \quad |w\rangle = |e_0\rangle + |e_1\rangle.
	\end{equation}
	Substituting these representations into the inner product, we obtain
	\begin{equation}
		\begin{split}
			\langle v|w\rangle &= (|e_0\rangle + i|e_1\rangle)^\dag (|e_0\rangle + |e_1\rangle) \\
			&= (\langle e_0| - i\langle e_1|)(|e_0\rangle + |e_1\rangle)  \\
			&= \langle e_0|e_0\rangle + \langle e_0|e_1\rangle -i \langle e_1|e_0\rangle -i\langle e_1|e_1\rangle \\
			&= 1-i.
		\end{split}
	\end{equation}
	We could equally well have obtained this result by simply substituting the definitions of $|v\rangle$ and $|w\rangle$ as column vectors into the inner product
	\begin{equation}
		\langle v|w\rangle \equiv \begin{pmatrix}
			1 & -i	
		\end{pmatrix} \begin{pmatrix}
			1 \\ 1
		\end{pmatrix} = 1-i.
	\end{equation}
\end{example}

Let us be honest: in two dimensions the first route is the more laborious one. From five or six dimensions onwards the balance tips.

\subsection{Inner Products}
Before we present the definition of a Hilbert space, we must explain the formal definition and a few properties of general inner products.

There are infinitely many inner products, and a vector space can perfectly well exist without one. We need one because quantum mechanics is probabilistic: it is our tool for probabilities.

An \emph{inner product} for a complex vector space $V$ assigns to any two vectors $|\phi\rangle, |\psi\rangle \in V$ a complex number $\langle \phi|\psi\rangle \in \mathbb{C}$. By definition it is additive in each argument (e.g., $\langle \phi_1 + \phi_2|\psi\rangle = \langle \phi_1|\psi\rangle + \phi_2|\psi\rangle$). For pulling out scalar factors $z\in \mathbb{C}$ the rules are
\begin{align}
	\langle \phi|z\psi\rangle &= z \langle \phi|\psi\rangle \\
	\langle z\phi|\psi\rangle &= z^* \langle \phi|\psi\rangle.
\end{align}
The complex conjugation in the second formula is important if the norm is to come out positive (see below).  

Part of the definition of the inner product is \emph{positivity}, in the sense that
\begin{equation}
	\begin{split}
		\langle \psi | \psi\rangle &\ge 0, \quad \forall \psi \in V.\\
		\langle \psi | \psi\rangle &= 0, \quad \text{only for $|\psi\rangle =0$}.\\
	\end{split}
\end{equation}
This property allows us to define the \emph{norm} of a vector $|\psi\rangle \in V$, often written $\|\psi\|$ or $\| |\psi\rangle \|$, as
\begin{equation}
	\| |\psi\rangle \| \equiv \sqrt{\langle \psi | \psi\rangle}.
\end{equation}
Two important consequences of positivity are
\begin{equation}
	\langle \phi | \psi\rangle = \langle \psi | \phi\rangle^*
\end{equation}
and the \emph{Schwarz inequality}
\begin{equation}
	|\langle \phi | \psi\rangle| \le \| \phi \|\|\psi \|.
\end{equation} 
The triangle inequality follows from this as well
\begin{equation}
	\left\|\sum_{j} |\psi_j\rangle \right\| \le \sum_{j} \|\psi_j\|.
\end{equation}
Both are working tools: the Schwarz inequality lies behind Heisenberg's uncertainty relation (Lecture 9), and we use the triangle inequality to bound quantities.

Two vectors are called \emph{orthogonal} (in symbols $|\phi\rangle \perp |\psi\rangle$) if
\begin{equation}
	\langle\phi|\psi\rangle = 0.
\end{equation}
Careful: a vanishing inner product does \emph{not} mean that one of the vectors is zero -- only the case $|\phi\rangle = |\psi\rangle$ forces that.


\subsection{Hilbert Space}


\begin{definition}
	A Hilbert space $\mathcal{H}$ is defined as a complex vector space with an inner product $\langle \cdot | \cdot  \rangle$ which is complete\footnote{The subtle aspect of completeness will play no great role in our lectures.} with respect to the norm $\|\psi\|^2 \equiv \langle \psi|\psi\rangle$.   
\end{definition}

For us, then, a Hilbert space is simply a complex vector space with an inner product.


\subsection{Bases}
We need bases for a practical reason: most calculations in quantum mechanics only become explicit once a basis has been fixed.

A \emph{basis} of a Hilbert space $\mathcal{H}$ is a system of vectors $|\phi_j\rangle \in \mathcal{H}$, where $j$ runs over some \emph{index set} $X$, with the properties
\begin{align}
	\langle \phi_j|\phi_k\rangle &=\delta_{jk}, \quad &\text{for all $j,k\in X$ (``$\{|\phi_j\rangle\}$ is \emph{orthonormal}'') and} \\
	\|\psi\|^2 &= \sum_{j\in X} |\langle \phi_j|\psi\rangle|^2, \quad &\text{for all $j,k\in X$ (``$\{|\phi_j\rangle\}$ is \emph{complete}'').}
\end{align}
Careful: this completeness has nothing to do with the completeness of the Hilbert space, and $\delta_{jk}$ is the Kronecker delta, not the Dirac delta function.

The most important property of such a basis is that every vector $|\psi\rangle \in \mathcal{H}$ can be expanded in it: we have
\begin{equation}
	|\Psi\rangle = \sum_{j\in X} \langle \psi_j|\Psi\rangle |\psi_j\rangle.
\end{equation}
This equation is one form of the so-called \emph{completeness relation}; written out with the basis vectors $|\phi_j\rangle$ it reads
\begin{equation}
	\eqbox{|\Psi\rangle = \sum_{j\in X} \langle \phi_j|\Psi\rangle\, |\phi_j\rangle.}
\end{equation}
It looks harmless, but it gets used all the time.

% fig_vl03_vollst.tex — Vollständigkeitsrelation: ein Vektor, zerlegt in
% zwei verschiedenen Orthonormalbasen desselben Raumes.
%
% Follows the conventions of fig_vl11_kasten.tex (the reference plate).
%
%   * geometry comes from the expansion derived in the lecture,
%         |Psi> = sum_j <phi_j|Psi> |phi_j>,
%     evaluated for the two orthonormal bases the lecture has already
%     introduced in C^2,
%         {|e_0>, |e_1>}   and   |e_±> = (|e_0> ± |e_1>)/sqrt(2),
%     on the concrete unit vector
%         |Psi> = (2|e_0> + |e_1>)/sqrt(5)   (real, so it fits in a plane).
%     Every coordinate below is that expansion, not a sketch:
%         c_0 = 2/sqrt5,  c_1 = 1/sqrt5   ->  |c|^2 = 4/5, 1/5
%         c_± = (2 ± 1)/sqrt10            ->  |c|^2 = 9/10, 1/10
%     and the two component rectangles therefore share |Psi> as their
%     diagonal exactly, because sum_j c_j |phi_j> = |Psi> in both bases.
%
%   * colour encodes *which basis*, and nothing else: cgblue for
%     {|e_0>,|e_1>}, birrenorange for {|e_+>,|e_->}.  The state |Psi>
%     itself is neutral navy — it belongs to no basis, which is the whole
%     point.  The coordinate rays are thin neutral grey: reference lines,
%     not data.
%
%   * within one basis the two members are told apart by tint (full
%     colour / 35%), not by a second hue, so the bars survive greyscale
%     printing; the numbers are printed as well, so nothing depends on
%     colour at all.  In (a) the same redundancy is carried by the dash
%     pattern — the first basis closes its rectangle with dashes, the
%     second with dots — and every arrow carries its own label, so the
%     two decompositions stay separable with the colour thrown away.
%     The neutral rays are held at birrengray!60: at full strength they
%     print at the same grey value as birrenorange and swallowed the
%     orange component vectors lying on top of them.
%
% Panel (a) the two decompositions of the same |Psi>.  Two different
%           rectangles, one diagonal.
% Panel (b) the weights |c_j|^2 those rectangles imply.  4/5 : 1/5 against
%           9/10 : 1/10 — the split is basis-dependent, the total is not.
%           This is the content of the completeness axiom
%           ||psi||^2 = sum_j |<phi_j|psi>|^2, and the reason the numbers
%           can later be read as probabilities.

\begin{center}
\begin{tikzpicture}[
    projvec/.style   ={-{Stealth[length=5pt,width=3.5pt]}, line width=1.3pt},
    statevec/.style  ={-{Stealth[length=6pt,width=4pt]}, line width=1.6pt,
                       spacecadet},
    basisray/.style  ={-{Stealth[length=4pt,width=3pt]}, line width=0.5pt,
                       birrengray!60},
    projdash/.style  ={line width=0.6pt, densely dashed},
    projdot/.style   ={line width=0.7pt, densely dotted},
]

% ─────────────────────────────────────────────────────────────
%  (a)  |Psi> resolved in two orthonormal bases
% ─────────────────────────────────────────────────────────────
\begin{axis}[
    name=geom,
    tpcplot,
    scale only axis,
    width=5.4cm, height=5.4cm,
    axis lines=none, grid=none,
    xmin=-0.38, xmax=1.38,
    ymin=-0.64, ymax=1.12,     % same span as x: 45 degrees looks like 45
    clip=false,
    title={(a) \tpcifenglish{One vector, two bases}%
                            {Ein Vektor, zwei Basen}},
]

  % ---- coordinate rays of the two bases (reference lines, neutral) ----
  \draw[basisray] (axis cs:-0.18,0)      -- (axis cs:1.24,0);
  \draw[basisray] (axis cs:0,-0.18)      -- (axis cs:0,1.00);
  \draw[basisray] (axis cs:-0.127,-0.127) -- (axis cs:0.721,0.721);
  \draw[basisray] (axis cs:-0.127,0.127)  -- (axis cs:0.424,-0.424);

  \node[lbl, anchor=west]       at (axis cs:1.27,0)      {$|e_0\rangle$};
  \node[lbl, anchor=south]      at (axis cs:0,1.03)      {$|e_1\rangle$};
  \node[lbl, anchor=south west] at (axis cs:0.74,0.74)   {$|e_+\rangle$};
  \node[lbl, anchor=north west] at (axis cs:0.44,-0.44)  {$|e_-\rangle$};

  % ---- the two component rectangles: same diagonal, different sides ----
  % basis {e_0,e_1}:  feet at (2/sqrt5, 0) and (0, 1/sqrt5)
  \draw[projdash, cgblue] (axis cs:0.894427,0)
                       -- (axis cs:0.894427,0.447214);
  \draw[projdash, cgblue] (axis cs:0,0.447214)
                       -- (axis cs:0.894427,0.447214);
  % basis {e_+,e_-}:  feet at 3/sqrt10 along e_+ and 1/sqrt10 along e_-
  \draw[projdot, birrenorange] (axis cs:0.670820,0.670820)
                            -- (axis cs:0.894427,0.447214);
  \draw[projdot, birrenorange] (axis cs:0.223607,-0.223607)
                            -- (axis cs:0.894427,0.447214);

  % ---- the component vectors c_j |phi_j> ----
  \draw[projvec, cgblue]       (axis cs:0,0) -- (axis cs:0.894427,0);
  \draw[projvec, cgblue]       (axis cs:0,0) -- (axis cs:0,0.447214);
  \draw[projvec, birrenorange] (axis cs:0,0) -- (axis cs:0.670820,0.670820);
  \draw[projvec, birrenorange] (axis cs:0,0) -- (axis cs:0.223607,-0.223607);

  % ---- the state itself ----
  \draw[statevec] (axis cs:0,0) -- (axis cs:0.894427,0.447214);
  \node[lbl, text=spacecadet, font=\small, anchor=south west]
    at (axis cs:0.91,0.47) {$|\Psi\rangle$};

  % ---- component labels, each beside its own arrow ----
  \node[lbl, text=cgblue, font=\scriptsize, anchor=north]
    at (axis cs:0.72,-0.045) {$c_0|e_0\rangle$};
  \node[lbl, text=cgblue, font=\scriptsize, anchor=east]
    at (axis cs:-0.04,0.22)  {$c_1|e_1\rangle$};
  \node[lbl, text=birrenorange, font=\scriptsize, anchor=south east]
    at (axis cs:0.58,0.685)  {$c_+|e_+\rangle$};
  \node[lbl, text=birrenorange, font=\scriptsize, anchor=north east]
    at (axis cs:0.155,-0.255) {$c_-|e_-\rangle$};

  \node[lbl, text=birrengray, font=\scriptsize, anchor=north east]
    at (axis cs:1.38,1.12) {$c_j \equiv \langle \phi_j|\Psi\rangle$};
\end{axis}

% ─────────────────────────────────────────────────────────────
%  (b)  The weights those two decompositions carry
% ─────────────────────────────────────────────────────────────
\begin{axis}[
    tpcplot,
    scale only axis,
    at={(geom.north east)}, anchor=north west, xshift=22mm,
    width=5.0cm, height=4.2cm,
    axis lines=left, grid=none,
    xmin=0, xmax=1.58, ymin=0.42, ymax=2.55,
    xtick={0,0.5,1}, xticklabels={$0$,$1/2$,$1$},
    ytick={1,2},
    yticklabels={{$\{|e_+\rangle,|e_-\rangle\}$},{$\{|e_0\rangle,|e_1\rangle\}$}},
    ytick style={draw=none},
    xlabel={$|c_j|^2$},
    clip=false,
    title={(b) \tpcifenglish{Weights}{Gewichte}},
]

  % the common total, drawn first so it does not sit on the bar borders
  \draw[levellinegrey] (axis cs:1,0.50) -- (axis cs:1,2.46);

  % basis {e_0,e_1}:  4/5 and 1/5
  \fill[cgblue]     (axis cs:0,1.74)  rectangle (axis cs:0.8,2.26);
  \fill[cgblue!35]  (axis cs:0.8,1.74) rectangle (axis cs:1.0,2.26);
  \draw[spacecadet, line width=0.5pt]
                    (axis cs:0,1.74)  rectangle (axis cs:1.0,2.26);

  % basis {e_+,e_-}:  9/10 and 1/10
  \fill[birrenorange]    (axis cs:0,0.74)  rectangle (axis cs:0.9,1.26);
  \fill[birrenorange!35] (axis cs:0.9,0.74) rectangle (axis cs:1.0,1.26);
  \draw[spacecadet, line width=0.5pt]
                         (axis cs:0,0.74) rectangle (axis cs:1.0,1.26);

  \node[font=\footnotesize, text=white] at (axis cs:0.40,2.00)
    {$|c_0|^2=\tfrac45$};
  \node[lbl, text=cgblue, font=\footnotesize, anchor=west]
    at (axis cs:1.04,2.00) {$|c_1|^2=\tfrac15$};

  \node[font=\footnotesize, text=white] at (axis cs:0.45,1.00)
    {$|c_+|^2=\tfrac{9}{10}$};
  \node[lbl, text=birrenorange, font=\footnotesize, anchor=west]
    at (axis cs:1.04,1.00) {$|c_-|^2=\tfrac{1}{10}$};

  \node[lbl, text=birrengray, font=\scriptsize]
    at (axis cs:0.50,1.50)
    {$\textstyle\sum_j |c_j|^2 = \|\Psi\|^2 = 1$};
\end{axis}

\end{tikzpicture}
\end{center}


The figure works this out for the two bases of the opening example,
$\{|e_0\rangle,|e_1\rangle\}$ and $\{|e_+\rangle,|e_-\rangle\}$.
In (a) each of the two decompositions of the same $|\Psi\rangle$ closes up into
a rectangle of its own -- different sides, the same diagonal. Panel (b) shows
what stays independent of the basis: not the individual weights, but their
sum.

For infinite-dimensional vector spaces the existence of bases is hard to prove. From now on we shall assume that we always have a basis; arbitrarily chosen vectors, after all, normally do \emph{not} form one. In practice bases are often obtained by the so-called orthonormalisation procedure: you start with some list of vectors $|\Phi_1\rangle$, $|\Phi_2\rangle$, $|\Phi_3\rangle$, $\ldots$. From these you determine, step by step, a system of pairwise orthogonal vectors $|\Psi_n\rangle$ according to the formula
\begin{equation}
	|\Psi_n\rangle \equiv |\Phi_n\rangle - \sum_{k=1}^{n-1} |\Psi_k\rangle \frac{\langle \Psi_k|\Phi_n\rangle}{\|\Psi_k\|^2},
\end{equation}
where summands with $|\Phi_k\rangle= 0$ are left out.\footnote{What is meant is $|\Psi_k\rangle = 0$: otherwise $\|\Psi_k\|^2 = 0$ would stand in the denominator and we would be dividing by zero.} This vector is orthogonal to the $|\Psi_k
\rangle$ with $k < n$. The vectors
\begin{equation}
	\left\{\frac{|\Psi_n\rangle}{\|\Psi_n\|}\middle| |\Psi_n\rangle\not=0\right\}
\end{equation}
then form an orthonormal system spanning the same subspace as the original list. If there is no vector perpendicular to all of the $|\Phi_n\rangle$, then this list is a basis as well.

Let us work the procedure through once in full.

\begin{example}
	We work in $\mathbb{C}^2$ and start with two made-up vectors
	\begin{equation}
		|\Phi_1\rangle \equiv \begin{pmatrix}
			1 \\ 0
		\end{pmatrix} \quad\text{and}\quad |\Phi_2\rangle \equiv \begin{pmatrix}
			1 \\ 1
		\end{pmatrix}.
	\end{equation}
	In the first step the sum runs from $1$ to $0$, so nothing is summed at all:
	\begin{equation}
		|\Psi_1\rangle = |\Phi_1\rangle = \begin{pmatrix}
			1 \\ 0
		\end{pmatrix}.
	\end{equation}
	In the second step only the term $k=1$ survives. With $\langle \Psi_1|\Phi_2\rangle = 1$ and $\|\Psi_1\|^2 = 1$ we obtain
	\begin{equation}
		|\Psi_2\rangle = |\Phi_2\rangle - |\Psi_1\rangle \frac{\langle \Psi_1|\Phi_2\rangle}{\|\Psi_1\|^2} = \begin{pmatrix}
			1 \\ 1
		\end{pmatrix} - \begin{pmatrix}
			1 \\ 0
		\end{pmatrix} = \begin{pmatrix}
			0 \\ 1
		\end{pmatrix}.
	\end{equation}
	Both vectors already have norm one: an arbitrary list has turned into the canonical basis $\{|e_0\rangle, |e_1\rangle\}$. In higher dimensions the calculation merely gets more laborious -- but work it does, every time.
\end{example}

\subsection{Operators}
Now to the second pillar of the theory. The vectors describe the \emph{states}, the operators the \emph{observables} -- and as soon as you have states and observables, you have a physical theory.

Let $\mathcal{H}$ be a Hilbert space. An \emph{operator} $A$ on $\mathcal{H}$ is a \emph{linear map}
\begin{equation}
	A:\mathcal{H}\rightarrow \mathcal{H}.
\end{equation}
That is to say,
\begin{equation}
	A(|\phi\rangle + |\psi\rangle) = A|\phi\rangle + A|\psi\rangle
\end{equation}
and
\begin{equation}
	A(\alpha|\phi\rangle) = \alpha A|\phi\rangle
\end{equation}
for all $\alpha\in \mathbb{C}$ and $|\phi\rangle, |\psi\rangle\in \mathcal{H}$.

\begin{example}
	In the case that $\mathcal{H}$ is finite-dimensional, operators can be represented by \emph{matrices}. For example, let us suppose that the dimension of $\mathcal{H}$ is $n=2$. All operators on $\mathcal{H}\equiv\mathbb{C}^2$ have the form of $2\times 2$ matrices:
	\begin{equation}
		A \equiv \begin{pmatrix}
			a_{11} & a_{12} \\ a_{21} & a_{22}
		\end{pmatrix},
	\end{equation} 
	where $a_{jk} \in \mathbb{C}$, $\forall 1 \le j,k\le 2$. The rules for multiplying matrices and vectors correspond exactly to the definition of linear operators.
\end{example}